\chapter{Математическая модель периодичности воспроизводства основного капитала}
\label{app:model}

\subsection{Периодичность в рамках одного предприятия}
Пусть предприятие $f_{j} \in \mathrm{F}_{N}$ каждый год $t \in \mathrm{T}$ производит и реализует товар $T_{t}$ со стоимостью $w_{t}$, равной:
\begin{equation}
    w_t = \underbrace{\overline{c}^o_t}_{\text{основной капитал}} + \underbrace{c^{wc}_t + v^{wc}_t}_{\text{оборотный капитал}} + \underbrace{m_t}_{\text{прибавочная стоимость}}, \quad \forall t \in \mathrm{T}.
\end{equation}
После реализации товара $T_{t} \rightarrow D_{t}$ годовая выручка $D_{t}$ распределяется следующим образом:
\begin{equation}D_t = \underbrace{\overline{c}^o_t}_{\text{фонд амортизации}} + \underbrace{c^{wc}_t}_{\text{возмещение сырья}} + \underbrace{v^{wc}_t}_{\text{выплата зарплаты}} + \underbrace{m_t}_{\text{прибыль}}, \quad \forall t \in \mathrm{T}.\end{equation}

Предположим, предприятие $f_{j} \in \mathrm{F}_{N}$ в период $t_{*} \in \mathrm{T}$ авансирует основной капитал в количестве $k_{*} \in \mathbb{N}_{0}$ единиц, стоимостном выражении $C^{o}_{*} \in \mathbb{R}_{\geq 0}$ и сроком физического износа $l\in \mathbb{N}$. При условии равномерного износа основного капитала среднегодовая стоимость его амортизации $\overline{c}^o_t$  в год $t \in \mathrm{T}$ составит
\begin{equation}\overline{c}^o_t = \frac{C^o_*}{l}, \quad \forall t \in \mathrm{T},\end{equation}
а остаточная стоимость основного капитала $C^{o}(\cdot)$ на конец периода $t \in \mathrm{T}$ будет равняться
\begin{equation}\begin{aligned} C^{o}(\cdot) : \{t_{*}, \dots, t_{*} +(l-1) \} \longrightarrow \mathbb{R}_{\geq 0} \\ t \longmapsto C^{0}(t) \end{aligned}, \quad C^{o}(t) = C^{0}_{*} - \underbrace{\sum_{n=t_{*}}^{t} \overline{c}^o_n}_{A} = C^{o}_{*} \cdot \underbrace{\left(1 - \frac{(t-t_{*})+1}{l} \right)}_{A},\end{equation}
где $A(\cdot)$ представляет собой накопленную сумму амортизации основного капитала на конец периода $t \in \mathrm{T}$:
$$\begin{aligned} A(\cdot) : \{t_{*}, \dots, t_{*} +(l-1) \} \longrightarrow \mathbb{R}_{\geq 0} \\ t \longmapsto A(t) \end{aligned}, \quad A(t) = \sum_{n=t_{*}}^{t} \overline{c}^o_n = C^{o}_{*} - C^{o}_{t}.$$
В таком случае на конец периода $t =[t_{*} + (l-1)] \in \mathrm{T}$ накопленная сумма амортизации $A(t)$ полностью покрывает первоначальную стоимость $C^{o}_{*}$ основного капитала:
$$A[t_{*} + (l-1)] = C^{o}_{*}.$$

Возмещение $\mathcal{I}(\cdot)$ основного капитала в объеме $k \in \mathbb{N}_{0}$ единиц осуществляется за счет средств накопленной амортизации $A_t = C^{o}_{*}$ в период $t \in \mathrm{T}$ в случае если в предыдущем периоде $(t-1) \in \mathrm{T}$ остаточная стоимость основного капитала $C^{o}_{t}$ равняется нулю:
\begin{equation}\begin{aligned} \mathcal{I}(\cdot) : \mathrm{T} \longrightarrow \{0,1\} \\ t \longmapsto \mathcal{I}(t) \end{aligned}, \quad \mathcal{I}(t) = \mathbf{1}_{C^{o}_{t-1}}(t) = \begin{cases} 1, &C^{o}_{t-1} = 0 \\ 0, &C^{o}_{t-1} \geq 0 \end{cases}.\end{equation}
При простом воспроизводстве количество авансируемых в периоде $t \in \mathrm{T}$ единиц $k \in \mathbb{N}_{0}$ можно представить в виде потока $k(\cdot)$:
\begin{equation}\begin{aligned}k(\cdot) : \mathrm{T} \longrightarrow \mathbb{N}_{0} \\ t \longmapsto k(t) \end{aligned}, \quad k(t) = k_{*} \mathbf{1}_{C^{o}_{t-1}}(t) = \begin{cases} k_*, &C^{o}_{t-1} = 0 \\ 0, &C^{o}_{t-1} \geq 0 \end{cases}.\end{equation}
Таким образом, периоды $t^{ce} \in \mathrm{T}^{ce} \subseteq \mathrm{T}$, характеризующиеся положительным значением потока $k(t) > 0$, представляют собой последовательность $\overrightarrow{\mathrm{T}}^{ce}$ с шагом $l \in \mathbb{N}$:
\begin{equation}\mathrm{T} \supseteq \mathrm{T}^{ce} = \{t^{ce} \mid k(t) > 0\} \implies \overrightarrow{\mathrm{T}}^{ce} = (t^{ce}_{i})_{i \in \mathrm{I}}, \quad t^{ce}_{i} = t_{*} + (i-1)l, \quad t^{ce}_{i+1} = t^{ce}_{i} + l,\end{equation}
а периоды $t^{da} \in \mathrm{T}^{da} \subseteq \mathrm{T}$ полного износа основного капитала $C^{o}(t) =0$ – последовательность $\overrightarrow{\mathrm{T}}^{da}$ с аналогичным шагом $l \in \mathbb{N}$:
\begin{equation}\mathrm{T} \supseteq \mathrm{T}^{da} = \{t^{da} \mid C^{o}(t)= 0\} \implies \overrightarrow{\mathrm{T}}^{da} = (t^{da}_{i})_{i \in \mathrm{I}}, \quad t^{da}_{i} = t_{*} + i(l-1), \quad t^{da}_{i+1} = t^{da}_{i} + l.\end{equation}
Составим общую последовательность $\overrightarrow{\mathrm{T}}^{o}$ из всех периодов $t^{o} \in \mathrm{T}^{o} \subseteq \mathrm{T}$, соответствующих полному износу $C^{o}(t)=0$ или полному возмещению $k(t)>0$ основного капитала:
\begin{equation}\mathrm{T} \supseteq \mathrm{T}^{o} = \{t^{ce}, t^{da} \} = \{t^{o}\} \implies \overrightarrow{\mathrm{T}}^{o} = (t^{o}_{i})_{i \in \mathrm{I}}, \quad t^{o}_{i} = \begin{cases} t_{2i-1} \hspace{-0.8em}&= t^{ce}_{i} \\ t_{2i} &= t^{da}_{i} \end{cases}, \quad \begin{cases} t^{ce}_{i} = t^{da}_{i-1} + 1 \\ t^{da}_{i} = t^{ce}_{i} + l \end{cases}.\end{equation}
В таком случае, процесс капиталистического воспроизводства в рамках одного предприятия можно разложить на непрерывный ежегодный оборот оборотного капитала $wc_{t} = c^{wc}_{t} + v^{wc}_{t}$ и дискретный оборот основного капитала $o=C^{o}_{*} = \sum_{n=t}^{t+(l-1)}\overline{c}^{o}_{n}$ в течение $l \in \mathbb{N}$ лет:
\begin{equation}\begin{cases} D^{wc}_t \rightarrow T^{wc}_t \rightarrow P_t^{wc}(c^{wc}_t + v^{wc}_t) \xRightarrow[\text{приращение стоимости}]{w_t^{wc} = c^{wc}_t + v^{wc}_t + m_t} T_t^{'wc} \rightarrow D_t^{'wc}, \quad D^{'wc}_t > D^{wc}_t \\ D^o_{t^{ce}_i} \rightarrow T^o_{t^{ce}_i} \rightarrow P^o_{t \in \{t^{ce}_{i}, \dots, t^{da}_{i}\}} \xRightarrow[\text{перенос стоимости}]{w^o_{t} = \overline{c}^{0}_{t}} T^o_{t^{da}_i} \rightarrow D^o_{t^{da}_i}, \quad D^{o}_{t^{ce}_{i}} = D^{o}_{t^{da}_{i}} \end{cases}.\end{equation}

\subsection{Периодичность в рамках выборки предприятий}
Пусть в экономике существует универсальное множество предприятий $\mathrm{F}_{\mathbb{U}}$. Определим $\mathrm{F}_{\mathbb{U}} \supseteq \mathrm F_N = \{f_j\}_{j \in \mathrm{J}^{N}}$ в качестве множества из $N \in \mathbb N$ предприятий, представляющего собой репрезентативную выборку, содержащую соответствующие критериям отбора $\mathcal{P}$ для анализа циклических процессов в капиталистической экономике предприятия $f_j \in \mathrm{F}_{N}$:
\begin{equation}\begin{aligned} \mathrm{F}_{\mathbb{U}} \supseteq \mathrm F_N := \{f \in \mathrm{F}_{\mathbb{U}} \mid \mathcal{P}(f)\} \\ |\mathrm{F}_{N}| = N, \; N \in \mathbb N \end{aligned} \implies \begin{aligned} \mathrm{F}_{N} = \{f_j \mid j \in \mathrm{J}^{N} \} \\ \mathrm{J}^{N} = \{ 1,2,\dots,N\} \end{aligned}.\end{equation}
В таком случае, базовый детерминированный поток обновления основного капитала $k^{\mathrm{d},j}(\cdot)$ на предприятии $f_j \in \mathrm{F}_{N}$ можно представить следующим образом:
\begin{equation}\begin{aligned} k^{\mathrm{d},j}(\cdot) : \mathrm{T} \longrightarrow \{0, k^{j}_{*} \} \\ t \longmapsto k^{\mathrm{d},j}(t) \\ k^{j}_{*} \in \mathbb{N} \end{aligned}, \quad k^{\mathrm{d},j}(t)=k^{j}_{*} \mathbf{1}_{\mathrm{T}^{\mathrm{d,ce},j}}(t) = \begin{cases} k^{j}_{*}, &t \in \mathrm{T}^{\mathrm{d,ce},j} \\ 0, &t \not\in \mathrm{T}^{\mathrm{d,ce},j} \end{cases},\end{equation}
где периоды обновления основного капитала $t^{\mathrm{d,ce},j} \in \mathrm{T}^{\mathrm{d,ce},j}$ образуют последовательность $\overrightarrow{\mathrm{T}}^{\mathrm{d,ce},j}$ с шагом $l^{j}$:
\begin{equation}\mathrm{T}^{\mathrm{d,ce},j} = \{t^{\mathrm{d,ce},j} \mid k^{\mathrm{d},j}(t^{\mathrm{d,ce},j}) = k^{j}_{*}\} \implies \begin{aligned} &\overrightarrow{\mathrm{T}}^{\mathrm{d,ce},j} = (t^{\mathrm{d,ce},j}_{i})_{i \in \mathrm{I}} \\ &t^{\mathrm{d,ce},j}_{i} = t^{j}_{*} + (i - 1)l^{j} \\ &t^{\mathrm{d,ce},j}_{i+1} = t^{\mathrm{d,ce},j}_{i} + l^{j} \end{aligned},\end{equation}
а совокупный квази-детерминированный поток обновления основного капитала $K^{\mathrm{qd}}(\cdot)$ по выборке $\mathrm{F}_{\mathbb{U}} \supseteq \mathrm F_N = \{f_j\}_{j \in \mathrm{J}^{N}}$:
\begin{equation}\begin{aligned} K^{\mathrm{qd}}(\cdot) : \mathrm{T} \longrightarrow \mathbb{N}_{0} \\ t \longmapsto K^{\mathrm{qd}}(t) \end{aligned}, \quad K^{\mathrm{qd}}(t) = \sum_{j \in \mathrm{J}^{N}} k^{\mathrm{d},j}(t) = \sum_{j \in \mathrm{J}^{N}} k^{j}_{*} \mathbf{1}_{\mathrm{T}^{\mathrm{d,ce},j}}(t),\end{equation}
со множеством периодов обновления $t^{\mathrm{qd,ce}} \in \mathrm{T}^{\mathrm{qd,ce}}$
\begin{equation}\mathrm{T}^{\mathrm{qd,ce}} = \{t^{\mathrm{qd,ce}} \mid K^{\mathrm{qd}}(t^{\mathrm{qd,ce}}) > 0\} \implies \mathrm{T}^{\mathrm{qd,ce}} = \{t^{\mathrm{d,ce},j} \mid t^{\mathrm{d,ce},j} \in \bigcup_{j \in \mathrm{J}^{N}} \mathrm{T}^{\mathrm{d,ce},j}\}.\end{equation}
Обозначим $\lambda(\cdot)$ в качестве определяющей уровень технологии производства на предприятии $f \in \mathrm{F}_{\mathbb{U}}$ функции. В таком случае, в силу однородности распространения научно-технического прогресса средний уровень технологии в рамках выборки $\mathrm{F}_{\mathbb{U}} \supseteq \mathrm F_N = \{f_j\}_{j \in \mathrm{J}^{N}}$ приблизительно совпадает для всех предприятий $f_j \in \mathrm{F}_{N}$, представляя собой поперечное среднее:
\begin{equation}\begin{aligned} \lambda(\cdot) : \mathrm{F}_{\mathbb{U}} \longrightarrow \mathbb{R}_{+} \\ f \longmapsto \lambda(f) \end{aligned}, \quad \begin{aligned} \mathbb{E}[\lambda(f_{\tilde{j}}) \mid \mathrm{F}_{N}] = \frac{1}{N} \sum_{j \in \mathrm{J}^{N}} \lambda_{j} =: \overline{\lambda}_{N} \\ \tilde{j} \sim \mathrm{unif}(\mathrm{J}^{N}) \end{aligned}, \quad \begin{aligned} \overline{\lambda}_{N} \longrightarrow \Lambda(\mathrm{F}_{N}) = \mathrm{const} \\ \max_{j\in \mathrm{J}^{N}}\,|\lambda_j-\Lambda|\xrightarrow[N\to\infty]{}0 \end{aligned}.\end{equation}
Срок физического износа основного капитала можно представить в качестве функции $l(\cdot, \cdot)$ от уровня технологии $\lambda(f)$ на предприятии $f \in \mathrm{F}_{\mathbb{U}}$, а также вектора прочих факторов $\hat{\hbar} = (\hbar_1, \dots, \hbar_n) \in \mathcal{H}^{n}$, оказывающих незначительное влияние:
\begin{equation}\begin{aligned} l(\cdot, \cdot) : \mathbb{R}_{+} \times \mathcal{H}^{n} \longrightarrow \mathbb{R}_{+} \\ (\lambda, \hat{\hbar}) \longmapsto l(\lambda, \hat{\hbar}) \end{aligned}, \quad \left| \frac{\partial l}{\partial \lambda} (\lambda, \hat{\hbar}) \right| \gg \|\nabla_{\hat{\hbar}} l(\lambda, \hat{\hbar})\|.\end{equation}
Следовательно, при фиксированных второстепенных факторах $\hbar = \mathrm{const}$ и сопоставимом уровне технологии производства $\Lambda = \mathrm{const}$ в рамках предприятий выборки $\mathrm{F}_{\mathbb{U}} \supseteq \mathrm F_N = \{f_j\}_{j \in \mathrm{J}^{N}}$ сроки физического износа основного капитала в среднем совпадают друг с другом:
\begin{equation}\begin{aligned} \hbar_j = \hbar = \mathrm{const} \quad \forall j \in \mathrm{J}^{N}\\ l_{\hbar} : \mathbb{R}_{+} \longrightarrow \mathbb{R}_{+} \\ \lambda \longmapsto l(\lambda, \hbar)  \end{aligned}, \quad \begin{aligned} \mathcal{G} = \sigma(\mathrm{F}_{N}, \Lambda, \hbar)\\ \mathbb E[l^{\tilde{j}}\mid\mathcal G]=L(\mathcal G) = \mathrm{const}\\ \tilde{j} \sim \mathrm{unif}(\mathrm{J}^{N})\end{aligned}, \quad \max_{j\in \mathrm{J}^{N}}\,|l^{j}-L|\xrightarrow[N\to\infty]{}0.\end{equation}
Положим в целях упрощения индивидуальные сроки физического износа основного капитала $l^{j}$ равными $L$ для всех предприятий $f_j \in \mathrm{F}_{N}$ в рамках выборки $\mathrm{F}_{\mathbb{U}} \supseteq \mathrm F_N = \{f_j\}_{j \in \mathrm{J}^{N}}$. В таком случае при простом воспроизводстве значение индивидуальных детерминированных потоков $k^{\mathrm{d},j}(t)$ после периода первоначального авансирования основного капитала $t^{j}_{*} \in \mathrm{T}$ повторяются для любых конгруэнтных по модулю срока физического износа $\mathrm{mod} \; L$ периодов $t \in \mathrm{T}$:
\begin{equation}\begin{aligned} l^{j} = L \\ \forall j \in \mathrm{J}^{N} \end{aligned} \implies \begin{aligned} \mathbf{1}_{\mathrm{T}^{\mathrm{d,ce},j}}(t_{1}) = \mathbf{1}_{\mathrm{T}^{\mathrm{d,ce},j}}(t_{2}) \\ \forall t_{1}, t_{2} \geq t^{j}_{*} : t_{1} \equiv t_{2} \; (\mathrm{mod} \; L) \end{aligned} \implies k^{\mathrm{d},j}(t_{1}) = k^{\mathrm{d},j}(t_{2}),\end{equation}
а следовательно, значения совокупного квази-детерминированного потока $K^{\mathrm{qd}}(t)$ также будут совпадать через каждые $L$ лет по истечении всех начальных периодов полного авансирования основного капитала $t^{j}_{*} \in \mathrm{T}^{*}$ в рамках предприятий выборки $\mathrm{F}_{\mathbb{U}} \supseteq \mathrm F_N = \{f_j\}_{j \in \mathrm{J}^{N}}$:
\begin{equation}K^{\mathrm{qd}}(t_{1}) = K^{\mathrm{qd}}(t_{2}) \quad \forall t_{1}, t_{2} \geq \dagger^{*} : t_{1} \equiv t_{2} \; (\mathrm{mod} \; L), \quad \dagger^{*} = \max_{j \in \mathrm{J}^{N}}t^{j}_{*}.\end{equation}
С течением времени под влиянием конкуренции квази-детерминированные периоды обновления основного капитала $t^{\mathrm{qd,ce}} \in \mathrm{T}^{\mathrm{qd,ce}}$ на предприятиях выборки $\mathrm{F}_{\mathbb{U}} \supseteq \mathrm F_N = \{f_j\}_{j \in \mathrm{J}^{N}}$ начинают синхронизироваться. Определим $\mathrm{T}^{\mathrm{d,ce}} \subseteq \mathrm{T}$ в качестве базовых детерминированных периодов обновления $t^{\mathrm{d,ce},} \in \mathrm{T}^{\mathrm{d,ce}}$ в рамках предприятий выборки $\mathrm{F}_{\mathbb{U}} \supseteq \mathrm F_N = \{f_j\}_{j \in \mathrm{J}^{N}}$ с промежутком в $L$ лет, образующих последовательность $\overrightarrow{\mathrm{T}}^{\mathrm{d,ce}}$:
\begin{equation}\mathrm{T}^{\mathrm{d,ce}} = \{t^{\mathrm{d,ce}}_{i}\}_{i \in \mathrm{I}} \implies \begin{aligned} &\overrightarrow{\mathrm{T}}^{\mathrm{d,ce}} = (t^{\mathrm{d,ce}}_{i})_{i \in \mathrm{I}} \\ &t^{\mathrm{d,ce}}_{i} = t_{*} + (i-1)L \\ &t^{\mathrm{d,ce}}_{i+1} = t^{\mathrm{d,ce}}_{i} + L \end{aligned}.\end{equation}
В таком случае при больших значениях $t \to \infty$ различия между квази-детерминированными периодами $t^{\mathrm{qd,ce}} \in \mathrm{T}^{\mathrm{qd,ce}}$ и базовыми детерминированными периодами $t^{\mathrm{d,ce}} \in \mathrm{T}^{\mathrm{d,ce}}$ устраняются, а множество квази-детерминированных периодов обновления основного капитала $\mathrm{T}^{\mathrm{qd,ce}}$ сводится ко множеству базовых детерминированных периодов $\mathrm{T}^{\mathrm{d,ce}}$ после некоторого порога $\dagger^{\mathrm{d}}$:
\begin{equation}\mathrm{dist}(t^{\mathrm{qd,ce}}, \mathrm{T}^{\mathrm{d,ce}}) = \min_{t^{\mathrm{d,ce}} \in \mathrm{T}^{\mathrm{d,ce}}} |t^{\mathrm{qd,ce}} - t^{\mathrm{d,ce}}| \xrightarrow[t^{\mathrm{qd,ce}} \to \infty]{t^{\mathrm{qd,ce}} \in \mathrm{T}^{\mathrm{qd,ce}}}0 \implies \exists \dagger^{\mathrm{d}} \in \mathrm{T} :  \mathrm{T}^{\mathrm{qd,ce}} \cap [\dagger^{\mathrm{d}}, \infty) \subseteq \mathrm{T}^{\mathrm{d,ce}}.\end{equation}
Определим базовый детерминированный поток обновления основного капитала $K^{\mathrm{d}}(\cdot)$, в рамках которого все предприятия выборки $\mathrm{F}_{\mathbb{U}} \supseteq \mathrm F_N = \{f_j\}_{j \in \mathrm{J}^{N}}$ синхронно обновляют основной капитал через каждые $L$ лет:
\begin{equation}\begin{aligned} K^{\mathrm{d}}(\cdot) : \mathrm{T} \longrightarrow \{0, K_{*} \} \\ t \longmapsto K^{\mathrm{d}}(t) \\ K_{*} \in \mathbb{N} \end{aligned}, \quad K^{\mathrm{d}}(t) = \mathbf{1}_{\mathrm{T}^{\mathrm{d,ce}}}(t) \underbrace{\sum_{j \in \mathrm{J}^{N}} k^{j}_{*}}_{K_{*}} = \begin{cases} K_{*}, &t \in \mathrm{T}^{\mathrm{d,ce}} \\ 0, &t \not\in \mathrm{T}^{\mathrm{d,ce}} \end{cases}.\end{equation}
В реальной капиталистической динамике фактические моменты обновления основного капитала могут отклоняться от базовых под воздействием конъюнктуры, кредита, конкуренции и прочих факторов. В таком случае, совокупный стохастический поток обновления основного капитала $K^{\mathrm{s}}(t)$ в рамках предприятий выборки можно определить как сумму базового детерминированного потока $K^{\mathrm{d}}(t)$ и совокупного отклонения $\Delta (\cdot, \cdot)$, зависящего от времени $t \in \mathrm{T}$ и вектора некоторых прочих факторов $\hat{z} = (z_{1}, \dots, z_{n}) \in \mathcal{Z}^{n}$:
\begin{equation}\begin{aligned} \Delta (\cdot, \cdot) : \mathrm{T} \times \mathcal{Z}^{n} \longrightarrow \mathbb{Z} \\ (t, \hat{z}) \longmapsto \Delta (t, \hat{z}) \end{aligned}, \quad K^{\mathrm{s}}(t) = K^{\mathrm{d}}(t) + \Delta (t, \hat{z}).\end{equation}
Обозначим $\xi_{t}$ в качестве текущих некоррелированных импульсов, представляющих собой исключительно количественные сдвиги относительно значений базового детерминированного потока замены $K^{\mathrm{d}}(t)$, не оказывающие влияния на совокупный запас основного капитала $K_{*} = \sum_{j \in \mathrm{J}^{N}} k^{j}_{*}$ предприятий выборки $\mathrm{F}_{\mathbb{U}} \supseteq \mathrm F_N = \{f_j\}_{j \in \mathrm{J}^{N}}$:
\begin{equation}\xi_{t} \sim \mathrm{unif} \left[-K^{\mathrm{d}}(t), \; K_{*} - K^{\mathrm{d}}(t) \right], \quad \begin{aligned} \mathrm{cov}(\xi_t,\xi_{t-n})=0 \\ \forall n\neq 0 \end{aligned}.\end{equation}
Поскольку текущие импульсы $\xi_{t}$ могут возникать исключительно в периоды $t^{\Phi^{4}} \in \mathrm{T}^{\Phi^{4}}$ фаз подъема $\Phi^{4} = \{\phi^{4}_{i}\}_{i \in \mathrm{I}}$ промышленного цикла (во время других фаз подавляющее большинство предприятий не обновляет основной капитал), определим индикатор $\mathcal{I}^{\xi}(\cdot)$:
\begin{equation}\begin{aligned} \mathcal{I}^{\xi}(\cdot) : \mathrm{T} \longrightarrow \{0,1\} \\ t \longmapsto \mathcal{I}^{\xi}(t) \end{aligned}, \quad \mathcal{I}^{\xi}(t) = \mathbf{1}_{\mathrm{T}^{\Phi^{4}}}(t) = \begin{cases} 1, &t \in \mathrm{T}^{\Phi^{4}} \\ 0, &t \not\in \mathrm{T}^{\Phi^{4}} \end{cases},\end{equation}
где $\phi(\cdot)$ представляет собой функцию фазировки промышленного цикла, сопоставляющую каждому периоду $t \in \mathrm{T}$ одну и четырех фаз промышленного цикла, на основе которой задается множество $\mathrm{T}^{\Phi^{4}}$ всех периодов фаз подъема $t^{\phi^{4}_{i}} \in \mathrm{T}^{\phi^{4}_{i}}$:
\begin{equation}\begin{aligned} \phi(\cdot) : \mathrm{T} \longrightarrow \{\phi^{1}, \phi^{2}, \phi^{3}, \phi^{4} \} \\ t \longmapsto \phi(t) \end{aligned}, \quad \mathrm{T}^{\phi^{4}_{i}} = \{t^{\phi^{4}_{i}} \mid \phi(t^{\phi^{4}_{i}}) = \phi^{4}_{i}\}, \quad \mathrm{T}^{\Phi^{4}} = \bigcup_{i \in \mathrm{I}} \mathrm{T}^{\phi^{4}_{i}}.\end{equation}
Таким образом, текущее отклонение $x(\cdot)$ в период $t \in \mathrm{T}$ представляет собой функцию
\begin{equation}\begin{aligned} x(\cdot) : \mathrm{T} \longrightarrow \left[-K^{\mathrm{d}}(t), \; K_{*} - K^{\mathrm{d}}(t) \right] \\ t \longmapsto x(t) \end{aligned}, \quad x(t) = \xi_{t} \, \mathbf{1}_{\mathrm{T}^{\Phi^{4}}}(t) = \begin{cases} \xi_{t}, &t \in \mathrm{T}^{\Phi^{4}} \\ 0, &t \not\in \mathrm{T}^{\Phi^{4}} \end{cases}.\end{equation}
При этом текущие отклонения $x_{t}$ в среднем взаимно компенсируют друг друга в рамках каждой из фаз подъема:
\begin{equation}\begin{aligned} K^{\mathrm{d}}(t) + x_{t} \leq K_{*} \\ \forall t \in \mathrm{T} \end{aligned}, \quad \begin{aligned} \mathrm{cov}(x_{t},x_{t-n})=0 \\ \forall n\neq 0 \end{aligned}, \quad \begin{aligned} \mathbb E\left[\sum_{t \in \mathrm{T}^{\phi^{4}_{i}}} x_{t} \mid \mathrm{T}^{\phi^{4}_{i}} \right] = 0 \\ \forall i \in \mathrm{I} \end{aligned}.\end{equation}
Поскольку часть отклонений имеет фазовый характер, возникая как сдвиг/разброс относительно базовых детерминированных периодов замены $t^{\mathrm{d,ce}} \in \mathrm{T}^{\mathrm{d,ce}}$ и изменяя тем самым фазу обновления основного капитала, в результате чего регулярно воспроизводится через кратные $L$ интервалы, обозначим $X(\cdot)$ в качестве накопленной фазовой деформации к периоду $t \in \mathrm{T}$:
\begin{equation}\begin{aligned} X(\cdot) : \mathrm{T} \longrightarrow \left[-K^{\mathrm{d}}(t), \; K_{*} - K^{\mathrm{d}}(t) \right] \\ t \longmapsto X(t), \end{aligned}, \quad X(t) = \sum_{n=1}^{\infty} x_{t-nL}, \quad X_{t} = X_{t-L} + x_{t-L}.\end{equation}
Чтобы учесть тенденцию капиталистической конкуренции к синхронизации (подтягивание к отраслевому ритму обновления, нивелирование рассогласований фаз), введем параметр $\rho\in[0,1)$, отражающий долю прошлой деформации, сохраняемую в следующем цикле обновления. В таком случае, накопленная фазовая деформация с затуханием $X^{\rho}(\cdot)$ представляет собой
\begin{equation}\begin{aligned} X^{\rho}(\cdot) : \mathrm{T} \longrightarrow \left[-K^{\mathrm{d}}(t), \; K_{*} - K^{\mathrm{d}}(t) \right] \\ t \longmapsto X^{\rho}(t), \end{aligned}, \quad X^{\rho}(t) = \sum_{n=1}^{\infty} \rho^{n} x_{t-nL}, \quad X_{t}^{\rho} = \rho (X^{\rho}_{t-L} + x_{t-L}),\end{equation}
а совокупное аддитивное отклонение с затуханием $\mathcal{X}^{\rho}(\cdot)$ в период $t \in \mathrm{T}$:
\begin{equation}\begin{aligned} \mathcal{X}^{\rho}(\cdot) : \mathrm{T} \longrightarrow \left[-K^{\mathrm{d}}(t), \; K_{*} - K^{\mathrm{d}}(t) \right] \\ t \longmapsto \mathcal{X}^{\rho}(t), \end{aligned}, \quad \mathcal{X}^{\rho}(t) = \sum_{n=0}^{\infty} \rho^{n} x_{t-nL} = \underbrace{\sum_{n=1}^{\infty} \rho^{n} x_{t-nL}}_{X_{t}^{\rho}} + x_{t} = X_{t}^{\rho} + x_t.\end{equation}
Обозначим $R(\cdot)$ в качестве потока нейтрализации, отвечающего за частичное нивелирование прошлых отклонений посредством их переноса на пиковый период базового детерминированного потока:
\begin{equation}\begin{aligned} R(\cdot) : \mathrm{T} \longrightarrow \left[-K^{\mathrm{d}}(t), \; K_{*} - K^{\mathrm{d}}(t) \right] \\ t \longmapsto R(t) \end{aligned}, \quad R(t)=(1-\rho)\,X^{\rho}_{t-L}\,\mathbf 1_{\mathrm{T}^{\mathrm{d,ce}}}(t).\end{equation}
В таком случае, совокупное отклонение $\Delta (t, \hat{z})$ можно представить в качестве функции $\Delta(\cdot, \cdot, \cdot)$ от текущего импульса $x_{t}$, затухающей накопленной фазовой деформации $X_{t}^{\rho}$ и потока нейтрализации $R_{t}$:
\begin{equation}\begin{aligned} \Delta(\cdot, \cdot, \cdot) : \{x_{t}, X^{\rho}_{t}, R_{t}\}_{t \in \mathrm{T}} \longrightarrow \left[-K^{\mathrm{d}}(t), \; K_{*} - K^{\mathrm{d}}(t) \right] \\ (x_{t}, X^{\rho}_{t}, R_{t}) \longmapsto \Delta (x_{t}, X^{\rho}_{t}, R_{t}) \end{aligned}, \quad \Delta (x_{t}, X^{\rho}_{t}, R_{t}) = x_{t} + X^{\rho}_{t} - R_{t},\end{equation}
или в виде следующей функции $\Delta(\cdot)$ от периода $t \in \mathrm{T}$:
\begin{equation}\begin{aligned} \Delta(\cdot) : \mathrm{T} \longrightarrow \left[-K^{\mathrm{d}}(t), \; K_{*} - K^{\mathrm{d}}(t) \right] \\ t \longmapsto \Delta(t) \end{aligned}, \quad \Delta(t) = \underbrace{\sum_{n=0}^{\infty} \rho^{n} x_{t-nL}}_{\mathcal{X}_{t}^{\rho}} - \underbrace{\mathbf 1_{\mathrm{T}^{\mathrm{d,ce}}}(t) \cdot (1- \rho) \sum_{n=1}^{\infty} x_{t-nL}}_{R_{t}}.\end{equation}
Следовательно, совокупный стохастический поток обновления основного капитала $K^{\mathrm{s}}(t)$ можно представить как
\begin{equation}K^{\mathrm{s}}(t) = K^{\mathrm{d}}(t) + \underbrace{\sum_{n=0}^{\infty} \rho^{n} x_{t-nL}}_{\mathcal{X}_{t}^{\rho}} - \underbrace{\mathbf 1_{\mathrm{T}^{\mathrm{d,ce}}}(t) \cdot (1- \rho) \sum_{n=1}^{\infty} x_{t-nL}}_{R_{t}}.\end{equation}